\chapter{Tổng quan vấn đề}

\section{Lý do chọn đề tài}
Bệnh tiểu đường là một trong những bệnh mạn tính phổ biến và có xu hướng gia tăng.
Việc xây dựng mô hình hỗ trợ dự đoán nguy cơ dựa trên dữ liệu y sinh có ý nghĩa thực tiễn
trong sàng lọc sớm. Bộ dữ liệu Pima Indians Diabetes được sử dụng rộng rãi trong giáo dục
và nghiên cứu vì có cấu trúc gọn, dễ tái lập và phù hợp để trình bày đầy đủ quy trình
rút trích dữ liệu và tiền xử lý.

\section{Mục tiêu nghiên cứu}
\begin{itemize}[leftmargin=1.5em]
    \item Xây dựng quy trình phân tích và tiền xử lý dữ liệu Pima có thể tái lập.
    \item Đánh giá ảnh hưởng của các chiến lược xử lý missing (Median, KNN, Iterative/MICE)
    và chuẩn hóa (StandardScaler, MinMaxScaler).
    \item So sánh hiệu năng giữa hai mô hình baseline: Logistic Regression và Random Forest.
\end{itemize}

\section{Câu hỏi nghiên cứu}
\begin{enumerate}[leftmargin=1.8em]
    \item Chính sách \textit{zero-as-missing} có cải thiện tính hợp lý thống kê và kết quả mô hình hay không?
    \item Tổ hợp \textit{imputer + scaler + model} nào phù hợp nhất với dữ liệu Pima trong điều kiện baseline?
    \item Trade-off giữa precision và recall thay đổi như thế nào khi điều chỉnh threshold?
\end{enumerate}

\section{Đối tượng và phạm vi}
Đối tượng nghiên cứu là bộ dữ liệu \texttt{pima-indians-diabetes.csv} (768 mẫu, 8 thuộc tính đầu vào và 1 nhãn đầu ra).
Báo cáo giới hạn trong phạm vi xử lý dữ liệu, EDA và đánh giá baseline; chưa mở rộng sang
hyperparameter tuning toàn diện hoặc bộ dữ liệu ngoài.

\section{Phương pháp tổng quát}
Nhóm sử dụng Python notebook để triển khai toàn bộ quy trình: nạp dữ liệu, phân tích mô tả,
làm sạch dữ liệu, tạo pipeline tiền xử lý, chia tập train/validation/test và so sánh định lượng.

\section{Kết cấu báo cáo}
Báo cáo gồm 5 chương:
\begin{itemize}[leftmargin=1.5em]
    \item Chương 1: Tổng quan vấn đề.
    \item Chương 2: Cơ sở lý thuyết và các khái niệm liên quan.
    \item Chương 3: Dữ liệu và phương pháp đề xuất.
    \item Chương 4: Thực nghiệm và kết quả đạt được.
    \item Chương 5: Kết luận và khuyến nghị.
\end{itemize}
